% SPDX-License-Identifier: CC-BY-SA-4.0
% Part: Syntactic analysis
% Section: Expressiveness of algebraic languages
% Exercise: Regular grammars

\begingroup

\begin{exercice}[Révisions -- Grammaires rationnelles]\label{exo:parsing/rational/revision}

  \begin{question}
  \item Poser le système d'équations correspondant à la grammaire ci-dessous. Quel est le problème lorsqu'on tente de le résoudre ?

    $$G_4 \eqdef \left\langle\{a, b\}, \{S, A\}, S, \left\{
    \begin{array}{lll}
      S   &\rightarrow& a A \mid \varepsilon \\
      A   &\rightarrow& S b\\
    \end{array}\right\}\right\rangle$$
  \end{question}

\end{exercice}

\begin{correction}

  \begin{question}
  \item Le système d'équations est $\left\{\begin{array}{lll}
    L   &=& a L_A \mid \varepsilon \\
    L_A   &=& L b\\
  \end{array}\right.$.
    Si on essaye de résoudre le système, on obtient $L = a L b \mid \varepsilon$ qui ne peut par être résolue avec le lemme d'Arden.
    Le langage est $\{a^n b^n | n\in \mathbb{N}\}$ qui n'est pas rationnel. 
  \end{question}
  
\end{correction}

\endgroup
\endinput
